\scanpage{27}
% archive sheet number: 26
% @ text / top
i.e.
% @ display / top, box
\[
\boxed{\Phi'=\Phi+(pL-F_S^*(L))}
\]
% @ text / upper third
On voit de $=$, si $\omega=\omega_{M,\varphi}\in\Gamma(\widehat\Omega^1_{S/W}\otimes(T\otimes_{\mathbb Z_p}\mathcal O_S))$ donne la forme connexion $\nabla\in\Gamma(\widehat\Omega^1_{S/W}\otimes M)$ via $\varphi$,

% @ display / upper third, connection matrix
\[
\nabla_{M,\varphi}=\begin{pmatrix}0&\omega_{M,\varphi}\\0&0\end{pmatrix}
\]
% @ display / right-hand box
\[
\boxed{\omega=\omega_{M,\varphi}\in\Gamma(\widehat\Omega^1_{S/W}\otimes(T\otimes_{\mathbb Z_p}\mathcal O_S))}
\]
% @ text / below matrix
qu'on a -- , si $\omega'=\omega_{M,\varphi'}$, qu'on a
% @ display / upper middle, box
\[
\boxed{\omega'_{\deleted{\unsure{M,\varphi}}}=\omega-dL}
\]
% @ text / middle
Écrivons les conditions \deleted{\unsure{qui|portant}} sur \deleted{\unsure{un tel}} $\Phi=\Phi_{M,\varphi}$ et $\omega=\omega_{M,\varphi}$ qui assurent qu'on a bien un objet de $C$. Ce sont

% @ display / middle to lower middle, brace with annotations
\[
\left\{\begin{aligned}
&\boxed{d\omega=0} &&\text{exprime que la connexion $\nabla$}\\
&&&\text{définie par $\omega$ est à courbure nulle}\\
&\boxed{F_S^*(\omega)-p\omega=d\Phi} &&\text{exprime que $F_M$ est}\\
&&&\text{compatible à la connexion}\\
&\boxed{\Phi_0=0} &&\text{exprimant que $(F_M)_0$}\\
&&&\text{respecte \deleted{\unsure{la}} $\mathrm{Fil}_0$.}
\end{aligned}\right.
\]
